%%%%%%%%%%%%%%%%%%%%%%%%%%%%%%%%%%%%%%%%%%%%%%%%%%%%%%%%%%%%%%%%%%%%%%
% Lớp 10 | Thống kê | Các số đặc trưng đo xu thế trung tâm của mẫu số liệu
% Dạng 0D6?3-2: Số trung bình cộng
% Số câu: 50  (sinh tự động bởi tach_cau.py; từ nay thầy sửa trực tiếp file này)
%%%%%%%%%%%%%%%%%%%%%%%%%%%%%%%%%%%%%%%%%%%%%%%%%%%%%%%%%%%%%%%%%%%%%%

%%%%%%%%%%%%%%%% TRẮC NGHIỆM 4 PHƯƠNG ÁN %%%%%%%%%%%%%%%%

% nguồn: MaHoaNH_TracNghiem_Lop10
\begin{ex}%[0D6N3-2]%[ma:29400d0d]
	Điểm kiểm tra môn Toán của một nhóm gồm $7$ học sinh như  sau
	$$\begin{array}{llllllll}7 & 5 & 6 & 8 & 9 & 10 & 6
\end{array}$$
Điểm trung bình môn Toán  của 7 học sinh đó (kết quả làm tròn đến hàng phần mười) là
\choice 
	 {$7{,}5$}
	 {$6{,}8$}
	 {$7{,}1$}
	{\True  $7{,}3$}
 \loigiai{
	 Ta có $\overline{x}=\dfrac{7+5+6+8+9+10+6}{7} \approx 7{,}3$.
} 
\end{ex}

% nguồn: MaHoaNH_TracNghiem_Lop10
\begin{ex}%[0D6N3-2]%[ma:d322a3a4]
	Ba nhóm học sinh gồm $10$ người, $15$ người, $25$ người.  Cân nặng trung bình của mỗi nhóm lần lượt là $50$ kg, $38$ kg, $40$ kg. Khối lượng trung bình của ba nhóm học sinh đó là
	\choice 
		{$26$  kg}
		{\True $41{,}4$ kg} 
		{$42{,}4$ kg} 
		{$37$  kg}
	\loigiai{ 
		 Khối lượng trung bình của ba nhóm học sinh là
		 \allowdisplaybreaks
		\begin{eqnarray*}
			\dfrac{1}{10+15+25}\cdot \left(10\cdot 50+15\cdot 38+25\cdot 40\right)=41{,}4\, \text{kg}. 
		\end{eqnarray*}
	 }
\end{ex}

% nguồn: MaHoaNH_TracNghiem_Lop10
\begin{ex}%[0D6N3-2]%[ma:b5b5739d]
	Chiều dài (đơn vị  mét) của 5 con cá voi xanh trưởng thành được cho như sau:
	$$ 
	\begin{array}{lllll}
		25 &  22 & 31 & 33 & 27.
	\end{array}
	$$ 
	Tìm chiều dài trung bình của $5$ con  cá voi xanh trên.
	\choice 
		 {$27$\,m}
		{\True $27{,}6$\,m} 
		 {$23$\,m}
		 {$28{,}2$\,m}
	\loigiai{ 
	Số trung  bình của mẫu số liệu $\dfrac{25+22+31+33+27}{5}\approx 27{,}6$ m.}
\end{ex}

% nguồn: MaHoaNH_TracNghiem_Lop10
\begin{ex}%[0D6N3-2]%[ma:9701c1c8]
	Tuổi đời của $16$ công nhân trong xưởng  sản xuất được thống kê trong bảng sau
	\begin{center}
		\begin{tabular}{|c|c|c|c|c|c|c|c|}
			 \hline Tuổi & $25$ & $26$ & $27$ & $29$ & $30$ & $33$ & Cộng\\
			\hline Số người & $2$ &  $3$ & $4$ & $3$ & $3$ & $1$ & $16$\\
			 \hline
		\end{tabular}
	\end{center}
	Tìm số trung bình  $\overline {x}$ của mẫu số liệu trên
	\choice 
		{$28$} 
		 {$27{,}75$}
		{\True  $27{,}875$}
		 {$27$}
	\loigiai{ 
		 Mẫu số liệu có số trung bình là
		 \allowdisplaybreaks
		\begin{eqnarray*}
			\overline {x} =\dfrac{25\cdot 2+26\cdot 3+27\cdot 4+29\cdot 3+30\cdot 3+33\cdot 1}{16}=  27{,}875.
		\end{eqnarray*}
	} 
\end{ex}

% nguồn: MaHoaNH_TracNghiem_Lop10
\begin{ex}%[0D6N3-2]%[ma:454396d1]
	Điều tra về  số tiền mua đồ dùng học tập trong một tháng của $40$ học sinh, ta có mẫu số liệu như sau (đơn vị: nghìn đồng)
	\begin{center}
		\begin{tabular}{|l|c|c|c|c|c|c|c|}
			\hline Giá trị  $(x)$ & {$[10; 15)$} & {$[15; 20)$} & {$[20; 25)$} & {$[25; 30)$} & {$[30; 35)$} & {$[35; 40)$} & Cộng\\
			\hline Tần số $(n)$ & $2$ & $5$ & $15$ &  $8$ & $9$ & $1$ & $\mathrm{~N}=40$\\
			 \hline
		\end{tabular}
	\end{center}
	Số trung bình của mẫu số liệu là 
	\choice 
		 {$22{,}5$}
		 {\True $25$}
		{$25{,}5$} 
		 {$27$}
	\loigiai{ 
		Số trung bình cộng của mẫu số  liệu là
		\allowdisplaybreaks 
		\begin{eqnarray*}
			\overline {x}=\dfrac{12{,}5\cdot 2+17{,}5\cdot 5+22{,}5\cdot  15+27{,}5\cdot 8+32{,}5\cdot 9+37{,}5\cdot 1}{40}=25.
		\end{eqnarray*}
	} 
\end{ex}

% nguồn: MaHoaNH_TracNghiem_Lop10
\begin{ex}%[0D6N3-2]%[ma:198624c0]
	Cho bảng phân bố tần số sau: khối lượng  $20$ học sinh lớp $10A$
	\begin{center}
		\begin{tabular}{|c|c|}
			\hline Khối lượng (kg) & Tần số\\ 
			\hline $50$ & $4$\\ 
			\hline $51$  & $5$\\
			\hline  $52$ & $6$\\
			\hline $55$  & $3$\\
			\hline $56$  & $2$\\
			\hline 
		\end{tabular}
	\end{center}
	Số trung bình cộng $\overline {x}$ của bảng  số liệu đã cho là
	\choice 
		{$53$} 
		 {$52{,}8$}
		{\True  $52{,}2$}
		{$52$} 
	 \loigiai{
		Giá trị  trung bình $\overline {x}=\dfrac{50\cdot 4+51\cdot 5+52\cdot 6+55\cdot 3+56\cdot 2}{20}=52{,}2$.
	} 
\end{ex}

% nguồn: MaHoaNH_TracNghiem_Lop10
\begin{ex}%[0D6N3-2]%[ma:fd466c67]
	Kết quả điểm kiểm tra 45 phút môn Hóa học của $100$ em học sinh được trình  bày ở bảng sau
	\begin{center}
		\begin{tabular}{|l|c|c|c|c|c|c|c|c|c|}
			\hline Điểm & $3$ & $4$ & $5$ & $6$ & $7$ & $8$ & $9$ &  $10$ & Cộng\\
			\hline Tần số & $3$ & $5$  & $14$ & $14$ & $30$ & $22$ & $7$ & $5$ & $100$\\
			 \hline
		\end{tabular}
	\end{center}
	Số trung bình cộng của bảng phân bố tần  số nói trên là
	 \choice
		 {\True $6{,}82$}
		{$4$} 
		{$6{,}5$} 
		 {$7{,}22$}
	 \loigiai{
		Số  trung bình cộng của bảng phân bố tần số nói trên là
		 \allowdisplaybreaks
		\begin{eqnarray*}
			\overline {x}=\dfrac{3\cdot 3+4\cdot  5+5\cdot 14+6\cdot 14+7\cdot 30+8\cdot 22+9\cdot 7+10\cdot 5}{100}=6{,}82.
		\end{eqnarray*}
	 }
\end{ex}

% nguồn: MaHoaNH_TracNghiem_Lop10
\begin{ex}%[0D6N3-2]%[ma:223bb8c9]
	Tuổi đời của $16$ công nhân trong xưởng sản xuất được  thống kê trong bảng sau
	\begin{center}
		\begin{tabular}{|l|l|l|l|l|l|l|l|}
			\hline Tuổi & $25$ & $26$  & $27$ & $29$ & $30$ & $31$ & Cộng\\
			\hline Số người & $2$ & $3$ & $4$ & $3$ & $3$ &  $1$ & $16$\\
			 \hline
		\end{tabular}
	\end{center}
	Số trung bình $\overline {x}$ của  mẫu số liệu trên là
	 \choice
		{$28$} 
		{\True  $27{,}75$}
		{$27{,}875$} 
		{$27$} 
	 \loigiai{
		Số trung bình $\overline {x}$ của mẫu  số liệu trên là
		\allowdisplaybreaks 
		\begin{eqnarray*}
			 \dfrac{25\cdot 2+26\cdot 3+27\cdot 4+29\cdot 3+30\cdot 3+31\cdot 1}{16}=27{,}75.
		\end{eqnarray*}
	 }
\end{ex}

% nguồn: MaHoaNH_TracNghiem_Lop10
\begin{ex}%[0D6N3-2]%[ma:5a259fec]
	Điểm số của $100$ học  sinh tham dự kỳ thi học sinh giỏi toán ở tỉnh $A$ (thang điểm là $20$) được thống kê theo bảng sau
	\begin{center}
		\begin{tabular}{|c|c|c|c|c|c|c|c|c|c|c|c|}
			\hline 
			Điểm $(x)$ & $9$ & $10$ & $11$  & $12$ & $13$ & $14$ & $15$ & $16$ & $17$ & $18$ & $19$\\
			 \hline
			Tần số $(n)$ & $1$ & $1$ & $3$ & $5$ & $8$ & $13$ & $19$ & $24$ & $14$ & $10$  & $2$\\
			 \hline
		\end{tabular}
	\end{center}
	Trung bình cộng của  bảng số liệu trên là
	 \choice
		{$15$} 
		{$15{,}50$} 
		{$16$} 
		 {\True $15{,}23$}
	 \loigiai{
		Trung bình cộng của bảng số  liệu trên là
		\allowdisplaybreaks 
		\begin{eqnarray*}
			\overline {x}=\dfrac{9\cdot 1+10\cdot 1+11\cdot 3+12\cdot 5+13\cdot 8+14\cdot 13+15\cdot 19+16\cdot 24+17\cdot  14+18\cdot 10+19\cdot 2}{100}=15{,}23.
		\end{eqnarray*}
	 }
\end{ex}

% nguồn: MaHoaNH_TracNghiem_Lop10
\begin{ex}%[0D6N3-2]%[ma:4dc3caf7]
	Thống kê số cuốn sách mỗi bạn trong lớp đã đọc trong năm $2022$, Nga thu được kết quả như bảng bên dưới. Hỏi trong năm $2022$, trung bình mỗi bạn trong lớp của  Nga đọc được bao nhiêu cuốn sách?
	\begin{center}
		\begin{tikzpicture}[xscale=1.7,yscale=.7]
				\begin{scope}[shift={(-.5,.5)}]
					\fill[cyan!10] (-2,0)  rectangle (0,-2);
					\draw (5,-2) grid (0,0) rectangle (-2,-2) 
						 (0,-1)--(-2,-1);
				\end{scope}
				\path 
					(-2.4,0) node[right]{Số  cuốn sách}
					(-2.4,-1)  node[right]{Số bạn}
					 (0,0) node{1}
					(0,-1)  node{3}
					 (1,0) node{2}
					 (1,-1) node{5}
					 (2,0) node{3}
					(2,-1)  node{15}
					(3,0)  node{4}
					 (3,-1) node{10}
					(4,0) node{5} 
					 (4,-1) node{7};
			\end{tikzpicture}
	\end{center}
	 \choice
		{\True  $3{,}325$}
		 {$2{,}325$}
		{$4{,}325$} 
		 {$3{,}5$}
	\loigiai{ 
	Trung bình mỗi bạn đọc được số  cuốn sách là
	\[ 
	\dfrac{1\cdot 3 + 2\cdot 5  + 3 \cdot 15 + 4 \cdot 10 + 5 \cdot 7}{3+5+15+10+7}= 3{,}325.
	 \]
	 }
\end{ex}

% nguồn: Ok_TracNghiemLop11_Lan1
\begin{ex}%[0D6N3-2]%[ma:0332998d]
	Cho mẫu số liệu về số sản phẩm đạt chất lượng tốt của một tổ sản xuất được trong một tháng như  sau
	\begin{center}
		\begin{tabular}{|l|c|c|c|c|c|}
			 \hline
			Số sản  phẩm&$90$&$100$&$110$&$120$&$130$\\
			 \hline
			Số  ngày&$1$&$4$&$8$&$11$&$6$\\
			\hline 
		\end{tabular}
	\end{center}
	Số trung bình trong mẫu số liệu  trên là
	 \choice
		{$116{,}(3)$} 
		 {$112$}
		 {\True $115{,}(6)$}
		{$110$} 
	 \loigiai{
		Số trung  bình trong mẫu số liệu trên là
	\[\overline{x}=\dfrac{90\cdot  1+100\cdot 4+110\cdot 8+120\cdot 11+130\cdot 6}{1+4+8+11+6}\approx 115{,}(6).\]}
\end{ex}

% nguồn: MaHoaNH_TracNghiem_Lop10
\begin{ex}%[0D6H3-2]%[ma:473ffabf]
	Điểm trung bình thi học kỳ II môn Toán của một nhóm gồm $N$ học sinh  lớp 12A6 là $8{,}1$. Biết rằng tổng điểm môn toán của nhóm này là $72{,}9$. Tìm số học sinh của nhóm.
	\choice 
		 {$20$}
		{\True  $9$}
		{$8$} 
		 {$15$}
	\loigiai{ 
		Ta có giá  giá trị $N=\dfrac{72{,}9}{8{,}1}=9$ (học sinh).
	} 
\end{ex}

% nguồn: MaHoaNH_TracNghiem_Lop10
\begin{ex}%[0D6H3-2]%[ma:3ef7d861]
	Để đủ điều kiện được cấp chứng chỉ A của trung tâm tin học, học viên phải trải qua $6$ lần thi  trắc nghiệm, thang điểm mỗi lần là $100$ điểm, và phải đạt trung bình $70$ điểm trở lên. Qua $5$ lần thi bạn Vân đạt trung bình $67{,}5$ điểm. Hỏi trong lần kiểm tra cuối cùng Vân phải đạt ít nhất bao nhiêu điểm để được cấp chứng chỉ?
	 \choice
		{$70$} 
		{$80$} 
		{$90$} 
		{\True  $82{,}5$}
	\loigiai{ 
		Gọi $x$ là điểm số trong lần kiểm tra cuối mà Vân cần đạt để được  cấp chứng chỉ.\\
		Ta có số điểm qua  $5$ lần thi của Vân là $67{,}5\cdot 5=337{,}5$.\\
		Theo bài ra ta có $\dfrac{x+337{,}5}{6}\geq 70 \Leftrightarrow x \geq 70\cdot  6-337{,}5 \Leftrightarrow x \geq 82{,}5$.\\
	Vậy trong lần kiểm tra cuối cùng, Vân phải đạt ít nhất $82{,}5$ điểm để được cấp chứng chỉ.} 
\end{ex}

% nguồn: MaHoaNH_TracNghiem_Lop10
\begin{ex}%[0D6H3-2]%[ma:e7a0a41e]
	Lớp $10$A của một trường trung học phổ thông có điểm thi môn Toán trong học kỳ $1$ vừa qua được thống  kê trong bảng sau
	\begin{center}
		\begin{tabular}{|l|c|c|c|c|c|c|}
			\hline Điểm thi & $5$ & $6$ & $7$ & $8$  & $9$ & $10$ \\
			\hline Tần số & $2$ & $7$ & $9$  & $12$ & $8$ & $4$ \\
			\hline 
		\end{tabular}
	\end{center}
	Tính điểm trung bình cộng môn Toán của lớp $10$A (làm  tròn đến hàng phần mười).
	\choice 
		 {$7{,}6$}
		 {$7{,}9$}
		{\True $7{,}7$} 
		{$8{,}0$} 
	\loigiai{ 
	Ta có $\overline{x}=\dfrac{5\cdot  2+6\cdot 7+7\cdot 9+8\cdot 12+9\cdot 8+10\cdot 4}{42}\approx 7{,}7$.}
\end{ex}

% nguồn: MaHoaNH_TracNghiem_Lop10
\begin{ex}%[0D6H3-2]%[ma:e813b229]
	Thời gian chạy (giây) cự li  $50$ m của $6$ học sinh được ghi lại trong bảng dưới đây
	\begin{center}
		\begin{tabular}{|c|c|c|c|}
			\hline Thời gian (giây) & $8{,}1$ &  $8{,}4$ & $8{,}5$ \\
			\hline  Tần số & $1$ & $2$ & $3$ \\
			 \hline
		\end{tabular}
	\end{center}
	Số trung bình cộng của mẫu số liệu trên  là
	\choice 
		{\True  $8{,}4$ (giây)}
		{$8{,}5$ (giây)} 
		 {$8{,}7$ (giây)}
		{$8{,}6$ (giây)} 
	\loigiai{ 
		Số  trung bình cộng của mẫu số liệu là
		$$\overline{x}=\dfrac{8{,}1\cdot  1+8{,}4\cdot 2+8{,}5\cdot 3}{1+2+3}=8{,}4.$$
	} 
\end{ex}

% nguồn: Ok_TracNghiemLop11_Lan1
\begin{ex}%[0D6B3-2]%[ma:35d4a71a]
	Cho $\sin \alpha =\dfrac{1}{2}$.  Giá trị của $P=\cos 2\alpha+\sin \alpha$ bằng
	 \choice
		{\True  $1$}
		{$2$} 
		 {$\dfrac{1}{2}$}
		{$0$} 
	\loigiai{Ta  có \allowdisplaybreaks
		\begin{eqnarray*}
			 P&=&\cos 2\alpha +3\sin \alpha\\
			 &=&1-2\sin^2 \alpha+\sin \alpha \\
			&=&1-2\left(\dfrac{1}{2}\right)^2+\dfrac{1}{2}=1. 
		\end{eqnarray*}
	 }
\end{ex}

% nguồn: Ok_TracNghiemLop11_Lan1
\begin{ex}%[0D6B3-2]%[ma:5a34098a]
	Rút  gọn biểu thức $P= \dfrac{1- \sin^2x}{\sin 2x} ~(\sin 2x \ne 0)$ ta được
	\choice 
		{$P= \dfrac{1}{2} \tan  x$}
		{\True $P= \dfrac{1}{2} \cot  x$}
		{$P=2 \cot  x$}
		{$P=2  \tan x$}
	 \loigiai{
	Ta có $P= \dfrac{1- \sin^2x}{\sin 2x}= \dfrac{\cos^2x}{2 \sin x \cos x}= \dfrac{\cos x}{2 \sin x}=  \dfrac{\cot x}{2}$.}
\end{ex}

% nguồn: Ok_TracNghiemLop11_Lan1
\begin{ex}%[0D6B3-2]%[ma:0ae68186]
	Tính  giá trị của biểu thức $P=\left(1-2\cos 2\alpha\right)\left(2+3\cos 2\alpha\right)$ biết $\sin\alpha=\dfrac{2}{3}$ .
	\choice 
		{$\True P=\dfrac{49}{27}$} 
		{$  P=\dfrac{50}{27}$}
		 {$ P=\dfrac{48}{27}$}
		{$ P=\dfrac{47}{27}$} 
	\loigiai{ 
		Ta có  \allowdisplaybreaks
		\begin{eqnarray*}
			P&=&\left(1-2\cos  2\alpha\right)\left(2+3\cos 2\alpha\right)\\
			&=&\left[1-2\left(1-2\sin^2\alpha\right)\right]\left[2+3\left(1-2\sin^2\alpha\right)\right] \\ 
			&=&\left[1-2\left(1-2\cdot \dfrac{4}{9}\right)\right]\left[2+3\left(1-2\cdot  \dfrac{4}{9}\right)\right]=\dfrac{49}{27}.
		\end{eqnarray*}
	} 
\end{ex}

% nguồn: Ok_TracNghiemLop11_Lan1
\begin{ex}%[0D6B3-2]%[ma:e54230e4]
	Cho $\sin \alpha=\dfrac{3}{4}$.  Khi đó, $\cos 2\alpha$ bằng
	 \choice
		 {$\dfrac{1}{8}$} {$\dfrac{\sqrt{7}}{4}$} {$-\dfrac{\sqrt{7}}{4}$} {\True $-\dfrac{1}{8}$}
	 \loigiai{
		Ta có, $\cos  2\alpha=1-2\sin^2 \alpha=1-2\left(\dfrac{3}{4}\right)^2=-\dfrac{1}{8}$.
	} 
\end{ex}


%%%%%%%%%%%%%%%% TRẢ LỜI NGẮN %%%%%%%%%%%%%%%%

% nguồn: MaHoaNH_TraLoiNgan_Lop10
\begin{bt}%[0D6N3-2]%[ma:0e723348]
	Cho mẫu số liệu có bảng tần suất  như sau
	\begin{center}
		\begin{tabular}{|c|c|c|c|c|c|}
			\hline 
			Giá trị $x_i$ & $4$ & $5$ &  $6$ & $7$ & $8$\\
			 \hline
			Tần số tương đối $f_i$ & $0{,}1$ & $0{,}45$ &  $0{,}2$ & $0{,}1$ & $0{,}15$\\
			 \hline
		\end{tabular}
	\end{center}
	Ta có số trung bình của mẫu số  liệu là
	 \shortans{5,75}
	 \loigiai{
		Ta có số trung  bình của mẫu số liệu là
		 \allowdisplaybreaks
		\begin{eqnarray*}
			\overline  {x}=4\cdot 0{,}1+5\cdot 0{,}45+6\cdot 0{,}2+7\cdot 0{,}1+8\cdot 0{,}15=5{,}75.
	\end{eqnarray*}}
\end{bt}

% nguồn: MaHoaNH_TraLoiNgan_Lop10
\begin{bt}%[0D6H3-2]%[ma:3d695b65]
	Một nhóm $11$ học sinh tham gia một kì thi. Số điểm thi của $11$ học sinh đó được sắp xếp từ thấp đến cao theo thang điểm $100$ như sau $0;0;63;65;69;70;72;78;81;85;89$. Tìm điểm  số trung bình của nhóm $11$ học sinh này?
	 \shortans{61,1}
	\loigiai{ 
		Điểm trung bình  là $\overline {x}=\dfrac{0+0+63+\ldots+85+89}{11}=\dfrac{672}{11}\approx 61{,}1$.
	 }
\end{bt}

% nguồn: Ok_TraLoiNgan_Lop11_Lan1
\begin{ex}%[0D6B3-2]%[ma:3126104e]
	Nếu $\sin x+\cos x=\dfrac{1}{2}$ thì $\sin 2 x$ bằng 
	 \shortans{-0,8}
	\loigiai{Ta có: $\sin x+\cos x=\dfrac{1}{2} \Rightarrow(\sin x+\cos x)^{2}=\dfrac{1}{4}  \Leftrightarrow 1+\sin 2 x=\dfrac{1}{4} \Leftrightarrow \sin 2 x=-\dfrac{3}{4}$.
	} 
\end{ex}

% nguồn: Ok_TraLoiNgan_Lop11_Lan1
\begin{ex}[c]%[0D6B3-2]%[ma:b3dc8976]
	Biết $\sin^4x=a+b\cos  2x+c \cdot \cos 4x$, với $a,b,c \in \mathbb{Q}$. Tính tổng $S=a+b+c$.
	 \shortans{0}
	 \loigiai{
		Ta có  \allowdisplaybreaks
		\begin{eqnarray*}
			& \sin^4x & =\left(\dfrac{1-\cos 2x}{2}\right)^2=\dfrac{1}{4}\left(1-2\cos 2x+\cos^22x\right)\\ 
			&  & =\dfrac{1}{4}\left(1-2\cos 2x+\dfrac{1+\cos 4x}{2}\right) =\dfrac{3}{8}-\dfrac{1}{2}\cos 2x+\dfrac{1}{8}\cos 4x.
		\end{eqnarray*}
	Vậy  $a=\dfrac{3}{8}$, $b=-\dfrac{1}{2}$, $c=\dfrac{1}{8}$. Nên $S=a+b+c=0$.}
\end{ex}


%%%%%%%%%%%%%%%% TỰ LUẬN %%%%%%%%%%%%%%%%

% nguồn: MaHoaNH_TuLuan_Lop10
\begin{bt}%[0D6N3-2]%[ma:e2064df4]
	Một  công ty vận chuyển $A$ dự kiến thưởng cho nhân viên giao hàng $B$ vào cuối năm dựa vào số đơn hàng giao được trong năm. Số đơn hàng của nhân viên $B$ giao được trong các tháng được cho trong dãy sau
	\begin{center}
		$1002$; $510$; $430$; $395$; $400$; $401$; $396$; $299$; $450$; $450$; $560$;  $611$.
	\end{center}
	Tính số đơn hàng  trung bình giao được trong $1$ tháng của nhân viên $B$.
	 \loigiai{
		Số đơn hàng trung bình giao được trong $1$ tháng của nhân viên  $B$:
		\allowdisplaybreaks 
		\begin{eqnarray*}
			\overline{x} =\dfrac{1002+510+430+395+400+401+396+299+450+450+560+611}{12}=439 \quad \text{(đơn hàng)}. 
		\end{eqnarray*}
	 }
\end{bt}

% nguồn: Ok_TuanLuan_Toan11_Lan1
\begin{bt}%[0D6N3-2]%[ma:b8c76d5f]
	Bảng số liệu sau đây thống kê thời gian nảy mầm một loại hạt đậu mới  trong các điều kiện khác nhau
	\begin{center}
		\begin{tabular}{|c|c|c|c|c|c|c|}
			\hline 
			Thời gian (phút)& $420$ & $440$ & $450$ & $480$ & $500$ & $540$  \\
			\hline 
			Tần số& $8$ & $17$ & $18$ &  $16$ & $11$ & $10$ \\
			 \hline
		\end{tabular}
	\end{center}
	Tính giá trị trung bình $\overline{x}$ (làm tròn đến hai chữ số sau dấu phẩy) về thời gian nảy mầm loại hạt mới nói  trên.
	\loigiai{Áp dụng công thức tính số trung bình cho bảng  tần số ta có
		$$\overline{x}=\dfrac{420\cdot8+440\cdot17+450\cdot18+480\cdot16+500\cdot11+540\cdot10}{80}=469 \text{  (phút).}$$
	} 
\end{bt}

% nguồn: Ok_TuanLuan_Toan11_Lan1
\begin{bt}%[0D6N3-2]%[ma:fe3cfaaa]
	Khối lượng $30$  chi tiết máy được cho bởi bảng sau
	\begin{center}
		\begin{tabular}{|c|c|c|c|c|c|c|c|}
			\hline 
			Khối lượng (gam)&  $250$ & $300$ & $350$ & $400$ & $450$ & $500$ & Tổng \\
			 \hline
			Tần số& $4$ & $4$ &  $5$ & $6$ & $4$ & $7$ & $30$ \\
			 \hline
		\end{tabular}
	\end{center}
	Tính số trung bình $\overline{x}$ (làm tròn đến chữ số  thứ hai sau dấu phẩy) của bảng nói trên.
	\loigiai{Áp dụng công thức tính số trung bình cho bảng tần số ta có 
		$$\overline{x}=\dfrac{250\cdot4+300\cdot4+350\cdot5+400\cdot6+500\cdot7}{30}\approx 388{,}33 \text{ (gam).}$$ 
	} 
\end{bt}

% nguồn: Ok_TuanLuan_Toan11_Lan2
\begin{ex}%[0D6N3-2]%[ma:32e0af2a]
	Bảng số liệu sau đây thống kê thời gian nảy mầm một  loại hạt đậu mới trong các điều kiện khác nhau
	\begin{center}
		\begin{tabular}{|c|c|c|c|c|c|c|}
			\hline 
			Thời gian  (phút)& $420$ & $440$ & $450$ & $480$ & $500$ & $540$ \\
			 \hline
			Tần số& $8$ & $17$ &  $18$ & $16$ & $11$ & $10$ \\
			\hline 
		\end{tabular}
	\end{center}
	Tính giá trị trung bình $\overline{x}$ (làm tròn đến hai chữ số sau dấu phẩy) về  thời gian nảy mầm loại hạt mới nói trên.
	\loigiai{Áp dụng công thức tính số trung bình cho bảng tần  số ta có
		 $$\overline{x}=\dfrac{420\cdot8+440\cdot17+450\cdot18+480\cdot16+500\cdot11+540\cdot10}{80}=469 \text{ (phút).}$$
	 }
\end{ex}

% nguồn: Ok_TuanLuan_Toan11_Lan2
\begin{ex}%[0D6N3-2]%[ma:b226b614]
	Khối  lượng $30$ chi tiết máy được cho bởi bảng sau
	\begin{center}
		\begin{tabular}{|c|c|c|c|c|c|c|c|}
			 \hline
			Khối lượng (gam)& $250$ & $300$ & $350$  & $400$ & $450$ & $500$ & Tổng \\
			\hline 
			Tần số& $4$ & $4$  & $5$ & $6$ & $4$ & $7$ & $30$ \\
			 \hline
		\end{tabular}
	\end{center}
	Tính số trung bình $\overline{x}$ (làm tròn đến chữ số thứ hai sau dấu phẩy) của bảng nói trên. 
	\loigiai{Áp  dụng công thức tính số trung bình cho bảng tần số ta có
		$$\overline{x}=\dfrac{250\cdot4+300\cdot4+350\cdot5+400\cdot6+500\cdot7}{30}\approx 388{,}33  \text{ (gam).}$$
	} 
\end{ex}

% nguồn: MaHoaNH_TuLuan_Lop10
\begin{bt}%[0D6H3-2]%[ma:a841d31a]
	Một cửa hàng bán $6$ loại quạt điện với giá tiền là $100$; $150$; $300$; $350$; $400$; $500$ (nghìn đồng). Số quạt điện mà cửa  hàng bán ra trong mùa hè vừa qua được thống kê trong bảng tần số sau
	\begin{center}
		\begin{tabular}{|c|c|c|c|c|c|c|}
			\hline 
			Giá tiền & $100$ & $150$ & $300$ & $350$ & $400$ &  $500$\\
			\hline 
			Số quạt bán được & $15$ &  $25$ & $31$ & $26$ & $12$ & $4$\\
			 \hline
		\end{tabular}
	\end{center}
	Tìm số tiền trung bình mà cửa hàng thu  được khi bán mỗi chiếc quạt?
	 \loigiai{
		$\overline {x}=\dfrac{100\cdot 15+150\cdot  25+300\cdot 31+350\cdot 26+400\cdot 12+500\cdot 4}{15+25+31+26+12+4}\approx 269$ (nghìn).
	 }
\end{bt}

% nguồn: MaHoaNH_TuLuan_Lop10
\begin{bt}%[0D6H3-2]%[ma:ecf4a2ee]
	Cho bảng số liệu điểm  giữa học kì $1$ môn Toán của $43$ học sinh lớp $10$A như sau
	\begin{center}
		\renewcommand{\arraystretch}{1.4} 
		\begin{tabular}{|c|c|c|c|c|c|c|c|c|c|c|}
			 \hline
			Điểm & $3$ & $4$ & $5$ & $6$ & $7$ & $8$ &  $9$ & $10$& Cộng\\
			\hline 
			Tần số & $3$ & $6$ & $12$ & $8$ & $6$ & $5$ &  $2$ & $1$ & $43$\\
			\hline 
		\end{tabular}
	\end{center}
	Dựa vào bảng trên hãy tính điểm trung bình môn Toán của $43$ học sinh lớp $10$A (\textit{kết  quả làm tròn đến hàng phần chục}).
	\loigiai{ 
		Ta có 
		 $$\overline{x}=\dfrac{3\cdot 3+4\cdot 6+5\cdot 12+6\cdot 8+7\cdot 6+8\cdot 5+9\cdot 2+10\cdot 1}{3+6+12+8+6+5+2+1}\approx 5{,}8.$$
	 }
\end{bt}

% nguồn: Ok_TuanLuan_Toan11_Lan2
\begin{ex}%[0D6H3-2]%[ma:2a38ea8c]
	Anh Thắng ghi lại cự li $20$ lần ném tạ sắt  $3$ kg của mình ở bảng sau (đơn vị: mét):
	\begin{center}
		\begin{tabular}{|l|l|l|l|l|l|l|l|l|l|}
			\hline 12,1 & 12,9 & 10,2 &  10,9 & 12,2 & 11,5 & 12,5 & 9,3 & 12,3 & 9,7 \\
			\hline 12,3 & 11,5  & 11,2 & 9,8 & 11,5 & 12,1 & 9,5 & 11,5 & 10,9 & 13,1 \\
			\hline 
		\end{tabular}
	\end{center}
	\begin{enumerate}
		 \item Tính cự li trung bình của mỗi lần ném.
		\item Tổng hợp lại kết quả  ném tạ xa của anh Thắng vào bảng tần số ghép nhóm theo mẫu sau
		\begin{center}
			\begin{tabular}{|c|c|c|c|c|c|}
				\hline Cự li (m) & $[9{,}2; 10) $ & $[10; 10{,}8) $ & $[10{,}8; 11{,}6) $ & $[11{,}6; 12{,}4) $  & $[12{,}4; 13{,}2) $ \\
				\hline Số lần & $ ? $ & $ ? $ & $ ? $  & $ ? $ & $ ? $ \\
				\hline 
			\end{tabular}
		\end{center}
		\item Hãy  uớc lượng cự li trung bình mỗi lần ném tạ xa từ bảng tần số ghép nhóm trên.
		\item Khả năng anh Thắng ném tạ xa được khoảng bao nhiêu  mét là cao nhất?
	\end{enumerate}
	 \loigiai{
		\begin{enumerate}
			\item  $\bar{x}=\dfrac{12{,}1\cdot 2+12{,}9+10{,}2+\cdots+9{,}5+13{,}1}{20}=11{,}3535$.
			\item Bảng  tần số ghép nhóm:
			\begin{center}
				\begin{tabular}{|l|c|c|c|c|c|}
					\hline Cự li (m) & $[9{,}2;  10) $ & $[10; 10{,}8) $ & $[10{,}8; 11{,}6) $ & $[11{,}6; 12{,}4) $ & $[12{,}4; 13{,}2) $ \\
					 \hline Số lần & 4 & 1 & 7 & 5 & 3 \\
					 \hline
				\end{tabular}
			\end{center}
			\item  Bảng tần số theo giá trị đại diện:
			\begin{center}
				\begin{tabular}{|l|c|c|c|c|c|}
					\hline  Cự li $ (m) $ & $[9{,}2; 10) $ & $[10; 10{,}8) $ & $[10{,}8; 11{,}6) $ & $[11{,}6; 12{,}4) $ & $[12{,}4; 13{,}2) $ \\
					\hline Giá trị đại diện & 9,6 & 10,4 & 11,2 & 12  & 12,8 \\
					\hline Số lần & 4  & 1 & 7 & 5 & 3 \\
					\hline 
				\end{tabular}
			\end{center}
			\item  Cự li trung bình:
			$\bar{x}=\dfrac{9{,}6\cdot 4+10{,}4\cdot 1+11{,}2\cdot  7+12\cdot 5+12{,}8\cdot 3}{20}=11{,}28$.
			 \item Ta có: ${u_m}=10{,}8,{u_{m+1}}=11{,}6,{n_m}=7,{n_{m+1}}=5;{n_{m-1}}=1,{u_{m+1}}-{u_m}=0{,}8$.\\
			Do  đó: ${M_0}=10{,}8+\dfrac{7-5}{7-5+7-1}\cdot 0{,}8=11$.\\
			Khả năng anh Thắng ném tạ xa được khoảng $11$  mét là cao nhất.
		\end{enumerate}
	} 
\end{ex}

% nguồn: Ok_TuanLuan_Toan11_Lan2
\begin{ex}%[0D6H3-2]%[ma:828f058f]
	Điểm thi môn Toán của hai lớp 11A và 11B được cho trong bảng sau, biết mỗi lớp có $50$ học sinh. Hỏi lớp nào có điểm  thi trung bình cao hơn?
	\begin{center}
		\begin{tabular}{|l|c|c|c|c|c|}
			\hline Điểm thi & $[0; 2) $ & $[2; 4) $ &  $[4; 6) $ & $[6; 8) $ & $[8; 10) $ \\
			\hline Lớp  11A & 2 & 4 & 12 & 28 & 4 \\
			\hline  Lớp 11B & 4 & 10 & 17 & 14 & 5 \\
			 \hline
		\end{tabular}
	\end{center}
	\loigiai{ 
		Ta có bảng thống kê điểm thi môn Toán theo giá trị đại  diện.
		\begin{center}
			\begin{tabular}{|l|c|c|c|c|c|}
				\hline Điểm thi & $[0; 2) $ & $[2; 4) $ & $[4; 6) $ &  $[6; 8) $ & $[8; 10) $ \\
				\hline Điểm thi  đại diện & 1 & 3 & 5 & 7 & 9 \\
				\hline Lớp 11A & 2 & 4 & 12 & 28 &  4 \\
				\hline Lớp  11B & 4 & 10 & 17 & 14 & 5 \\
				\hline 
			\end{tabular}
		\end{center}
		Điểm thi trung bình của lớp 11A là $( 1\cdot 2+3\cdot 4+5\cdot 12+7\cdot 28+9\cdot 4): 50=6{,}12$.\\ 
		Điểm thi trung bình của lớp 11B là $( 1\cdot 4+3\cdot 10+5\cdot 17+7\cdot  14+9\cdot 5): 50=5{,}24$.\\
		Vậy điểm thi trung bình của lớp  11A cao hơn lớp 11B.
	 }
\end{ex}

% nguồn: Ok_TuanLuan_Toan11_Lan2
\begin{ex}%[0D6H3-2]%[ma:995da9dc]
	Một thư viện thống kê số lượng  sách được mượn mỗi ngày trong $3$ tháng ở bảng sau:
	\begin{center}
		\begin{tabular}{|c|c|c|c|c|c|c|c|}
			\hline Số sách & $[16; 20] $ & $[21; 25] $ &  $[26; 30] $ & $[31; 35] $ & $[36; 40] $ & $[41; 45] $ & $[46; 50] $ \\
			\hline Số ngày & 3 & 6 & 15 & 27 & 22 & 14  & 5 \\
			\hline 
		\end{tabular}
	\end{center}
	Trong  $3$ tháng trên số sách được mượn mỗi ngày cao nhất là bao nhiêu?
	\loigiai{ 
		Do số sách là số nguyên nên ta hiệu chỉnh lại  như sau:\\
		\begin{tabular}{|>{\raggedright\arraybackslash}m{1.6cm}|>{\centering\arraybackslash}m{1.6cm}|>{\centering\arraybackslash}m{1.6cm}|>{\centering\arraybackslash}m{1.6cm}|>{\centering\arraybackslash}m{1.6cm}|>{\centering\arraybackslash}m{1.6cm}|>{\centering\arraybackslash}m{1.6cm}|>{\centering\arraybackslash}m{1.6cm}|>{\centering\arraybackslash}m{1.6cm}|>{\centering\arraybackslash}m{1.6cm}|>{\centering\arraybackslash}m{1.6cm}|}
			\hline Số sách & $[15{,}5; 20{,}5) $ & $[20{,}5; 25{,}5) $ & $[25{,}5; 30{,}5) $ & $[30{,}5; 35{,}5) $ & $[35{,}5; 40{,}5] $  & $[40{,}5; 45{,}5] $ & $[45{,}5; 50{,}5) $ \\
			\hline Số ngày & 3 &  6 & 15 & 27 & 22 & 14 & 5 \\
			\hline 
		\end{tabular}\\
		Ta có nhóm  chứa mốt của mẫu số liệu trên là nhóm $[30{,}5;35{,}5)$.\\
		Do đó ${u_m}=30{,}5$, ${n_{m-1}}=15$, ${n_m}=27$, ${n_{m+1}}=22$, ${u_{m+1}}-{u_m}=35{,}5-30{,}5=5$.\\ 
		Mốt của mẫu số liệu ghép  nhóm là: ${M_0}=30{,}5+\dfrac{27-15}{(27-15)+(27-22)}\cdot 5=\dfrac{1157}{34}\approx 34$.
	} 
\end{ex}

% nguồn: Ok_TuanLuan_Toan11_Lan2
\begin{ex}%[0D6H3-2]%[ma:c7dfd76f]
	Cân lần lượt $40$ quả cam (đơn vị gram) ta được  kết quả sau (mẫu số liệu)
	\begin{center}
		\begin{tabular}{cccccccccccccccc}
			85&86&86&86&86&86&87&87&87&87&87&88&88&88&88&89\\ 
			 89&89&89&89&89&89&90&90&90&90&91&91&91&92&93\\
			 93&93&93&94&94&94&94&94&94
		\end{tabular}
	\end{center}
	\begin{enumerate}
		\item Tính khối lượng trung bình của $40$  quả cam.
		 \item Tổng hợp lại số liệu trên vào bảng tần số ghép nhóm theo mẫu sau:
		\begin{center}
			\begin{tabular}{|l|c|c|c|c|c|}
				\hline Khối lượng & $[85; 87) $ &  $[87; 89) $ & $[89; 91) $ & $[91; 93) $ & $[93; 95) $ \\
				\hline Số quả & $ ? $ & $ ? $  & $ ? $ & $ ? $ & $ ? $ \\
				\hline 
			\end{tabular}
		\end{center}
		\item Hãy ước  lượng trung bình khối lượng của $40$ quả cam từ bảng tần số ghép nhóm trên.
		\item Dự đoán khối lượng  quả cam xuất hiện nhiều nhất.
	\end{enumerate}
	 \loigiai{
		\begin{enumerate}
			 \item Khối lượng trung bình của $40$ quả cam là:\\
			$\overline{x}=\dfrac{85+86\cdot 5+87\cdot 5+88\cdot 4+89\cdot  7+90\cdot 4+91\cdot 3+92+93\cdot 4+94\cdot 6}{40}=89{,}65$ (g).
			\item Bảng tần  số ghép nhóm:
			\begin{center}
				\begin{tabular}{|l|c|c|c|c|c|}
					\hline Khối lượng & $[85; 87) $ & $[87;  89) $ & $[89; 91) $ & $[91; 93) $ & $[93; 95) $ \\
					 \hline Số quả & 6 & 9 & 11 & 4 & 10 \\
					\hline 
				\end{tabular}
			\end{center}
			\item Bảng tần số  ghép nhóm theo giá trị đại diện:
			\begin{center}
				\begin{tabular}{|l|c|c|c|c|c|}
					\hline Khối lượng & $[85; 87) $ & $[87;  89) $ & $[89; 91) $ & $[91; 93) $ & $[93; 95) $ \\
					\hline Giá trị đại diện & 86 & 88 & 90 &  92 & 94 \\
					\hline Số  quả & 6 & 9 & 11 & 4 & 10 \\
					 \hline
				\end{tabular}
			\end{center}
			Khối lượng trung  bình của 40 quả cam là
			\begin{center}
				$\dfrac{86\cdot 6+88\cdot 9+90\cdot 11+92\cdot  4+94\cdot 10}{40}=90{,}15$ (g).
			\end{center}
			\item Nhóm chứa mốt  của mẫu số liệu trên là $[89;91)$.\\
			Do đó: ${u_m}=89;{n_{m-1}}=9;{n_{m+1}}=4;{u_{m+1}}-{u_m}=91-89=2$.\\ 
			Mốt của mẫu số liệu ghép nhóm  là:
			\begin{center}
				${M_0}=89+\dfrac{11-9}{( 11-9)+( 11-4)}\cdot  2=\dfrac{805}{9}\approx 89{,}44$ (g).
			\end{center}
			Vậy khả năng quả  cam nặng $89{,}44$ (g) là cao nhất.
		\end{enumerate}
	} 
\end{ex}

% nguồn: Ok_TuanLuan_Toan11_Lan2
\begin{ex}%[0D6H3-2]%[ma:6e44b8b7]
	Điều tra về thu nhập của công nhân xí nghiệp X (đơn vị: nghìn  đồng/tháng), người ta ghi được bảng tần số ghép lớp sau đây:
	\begin{center}
		\begin{tabular}{|c|c|}
			\hline Tiền lương &  Số người \\
			\hline$ [800; 890] $ & 15  \\
			$ [900; 990] $ & 25  \\
			$ [1000; 1090] $ & 28 \\ 
			$ [1100; 1190] $ &  35 \\
			$ [1200; 1290] $ & 40  \\
			$ [1300; 1390] $ & 30  \\
			$ [1400; 1490] $  & 27 \\
			 \hline
		\end{tabular}
	\end{center}
	Hãy  ước lượng số trung bình và mốt của mẫu số liệu ghép nhóm trên.
	 \loigiai{
		Hiệu chỉnh  lại bảng số liệu ta có:
		\begin{center}
			\begin{tabular}{|>{\raggedright\arraybackslash}m{1.5cm}|>{\centering\arraybackslash}m{1.5cm}|>{\centering\arraybackslash}m{1.5cm}|>{\centering\arraybackslash}m{1.75cm}|>{\centering\arraybackslash}m{1.75cm}|>{\centering\arraybackslash}m{1.75cm}|>{\centering\arraybackslash}m{1.75cm}|>{\centering\arraybackslash}m{1.75cm}|}
				\hline Tiền lương & $[795; 895) $ & $[895; 995) $ & $[995;  1095) $ & $[1095; 1195) $ & $[1195; 1295) $ & $[1295; 1395) $ & $[1395; 1495) $ \\
				\hline Giá trị đại diện & 845 & 945 & 1045 & 1145 & 1245 & 1345 &  1445 \\
				\hline Số người & 15 &  25 & 28 & 35 & 40 & 30 & 27 \\
				\hline 
			\end{tabular}
		\end{center}
		Thu nhập trung bình của công nhân xấp xỉ  bằng:
		\begin{center}
			$\overline{x}=\dfrac{845\cdot  15+945\cdot 25+1045\cdot 28+1145\cdot 35+1245\cdot 40+1345\cdot 30+1445\cdot 27}{200}=1174$.
		\end{center}
		Nhóm chứa  mốt của mẫu số liệu là $[1195;1295)$.\\
		 Do đó ${u_m}=1195;{n_{m-1}}=35,{n_{m+1}}=30,{u_{m+1}}-{u_m}=1295-1195=100$.\\
		Mốt của mẫu số liệu ghép nhóm  là:
		\begin{center}
			${M_0}=1195+\dfrac{40-35}{( 40-35)+(  40-30)}\cdot 100=\dfrac{3685}{3}\approx 1228{,}33$.
		\end{center}
	} 
\end{ex}

% nguồn: Ok_TuanLuan_Toan11_Lan2
\begin{ex}%[0D6H3-2]%[ma:84370d51]
	Kết quả đo  khối lượng của $30$ củ khoai tây ở nông trường được biểu diễn ở biểu đồ dưới đây.
	\begin{center}
		\begin{tikzpicture}[scale=1,>=stealth,font=\footnotesize,line join=round,line cap=round]
				\draw[->]  (0,0)--(7,0);
				\draw[->] (0,0)--(0,5)  node[midway,sloped,above=12pt]{\small{\bf Số cây}};
				\foreach  \y in {1,2,3,4}
				\pgfmathsetmacro\giatri{int  (3*\y)}
				\draw[gray] (0,\y)  node[left]{\color{black}{\giatri}}--(6.5,\y);
				\foreach \x/\y  in {1/1,2/1.66,3/4,4/2,5/1}
				\draw[fill=blue!50] (\x,0) rectangle  +(1,\y);
				\draw[dashed] (2,1.66)--(0,1.66)  node[left]{$5$};
				\foreach \x/\y in  {1/70,2/80,3/90,4/100,5/110,6/120}
				\draw (\x,0) node[below]{$\y$}; 
			\end{tikzpicture}
	\end{center}
	Hãy ước lượng số trung  bình và mốt của mẫu số liệu ghép nhóm trên.
	\loigiai{ 
		\begin{center}
			\begin{tabular}{|c|c|c|c|c|c|}
				\hline Khối lượng & $[70; 80) $ & $[80; 90) $ & $[90; 100) $ & $[100; 110) $  & $[110; 120) $ \\
				\hline Giá trị đại diện &  75 & 85 & 95 & 105 & 115 \\
				\hline Số củ & 3 & 5 & 12 & 6  & 3 \\
				\hline 
			\end{tabular}
		\end{center}
		Khối lượng trung bình mỗi  củ khoai tây xấp xỉ bằng:
		\begin{center}
			 $\dfrac{75\cdot 3+85\cdot 5+95\cdot 12+105\cdot 6+115\cdot 3}{30}=\dfrac{553}{6}\approx 92{,}17$ (g).
		\end{center}
		Nhóm  chứa mốt của mẫu số liệu là $[90;100)$.\\
		Do  đó ${u_m}=90;{n_{m-1}}=5,{n_{m+1}}=6,{u_{m+1}}-{u_m}=100-90=10$.\\
		Mốt của mẫu số liệu  ghép nhóm là:
		\begin{center}
			${M_0}=90+\dfrac{12-5}{( 12-5)+( 12-6)}\cdot  10=\dfrac{1240}{13}\approx 95{,}38$ (g).
		\end{center}
	 }
\end{ex}

% nguồn: Ok_TuanLuan_Toan11_Lan2
\begin{ex}%[0D6H3-2]%[ma:408c9c38]
	Nhiệt độ trung bình của 12 tháng tại thành phố Hà Nội từ năm  1961 đến hết năm 1990 (30 năm) được cho trong bảng sau:
	\begin{center}
		\begin{tabular}{|l|c|c|c|c|c|}
			\hline Nhiệt độ $ \left(^{\circ} \mathrm{C}\right) $ & $[12; 14)  $ & $[14; 16) $ & $[16; 18) $ & $[18; 20) $ & $[20; 22) $ \\
			\hline Tần số & 1  & 3 & 12 & 9 & 5 \\
			\hline 
		\end{tabular}
	\end{center}
	Tính mốt của mẫu số  liệu trên.
	 \loigiai{
		Nhóm  chứa mốt của mẫu số liệu trên là nhóm $[16;18)$.\\
		Nên ta có ${u_m}=16,{n_{m-1}}=3,{n_m}=12,{n_{m+1}}=9,{u_{m+1}}-{u_m}=18-16=2  $.\\
		Mốt của mẫu số liệu trên là ${M_0}=16+\dfrac{12-3}{(  12-3)+( 12-9)}\cdot 2=17{,}5 $.
	} 
\end{ex}

% nguồn: Ok_TuanLuan_Toan11_Lan2
\begin{ex}%[0D6H3-2]%[ma:5b121ecb]
	Thời gian (phút) đi từ nhà đến trường của bạn A trong $35$ ngày được cho trong bảng sau: 
	\begin{center}
		\begin{tabular}{|c|c|c|c|c|c|}
			\hline Thời gian (phút) & $[19; 21) $ & $[21; 23) $ & $[23; 25) $ & $[25; 27) $ & $[27; 29)  $ \\
			\hline Tần số & 5 & 8 & 7 & 9 & 6 \\ 
			 \hline
		\end{tabular}
	\end{center}
	Dự đoán thời gian bạn  A thường dành ra để đi từ nhà đến trường.
	 \loigiai{
		 Ta có bảng thống kê thời gian đi từ nhà đến trường của A theo giá trị đại diện.
		\begin{center}
			\begin{tabular}{|c|c|c|c|c|c|}
				\hline Thời gian (phút) & $[19;  21) $ & $[21; 23) $ & $[23; 25) $ & $[25; 27) $ & $[27; 29) $ \\
				\hline Thời gian đại diện & 20 & 22 & 24 & 26 & 28  \\
				\hline Tần sồ & 5 &  8 & 7 & 9 & 6 \\
				\hline 
			\end{tabular}
		\end{center}
		 Thời gian trung bình bạn A thường dành ra để đi từ nhà đến trường là
		\begin{center}
			$(  20\cdot 5+22\cdot 8+24\cdot 7+26\cdot 9+28\cdot 6): 35=24{,}2 $.
		\end{center}
	} 
\end{ex}

% nguồn: Ok_TuanLuan_Toan11_Lan2
\begin{ex}%[0D6H3-2]%[ma:98398a0b]
	Tuổi thọ (năm) của $50$ bình ắc quy  ô tô được cho như sau:
	\begin{center}
		\begin{tabular}{|l|c|c|c|c|c|c|}
			\hline Tuổi thọ (năm) & $[2; 2,5) $ &  $[2,5; 3) $ & $[3; 3,5) $ & $[3,5; 4) $ & $[4; 4,5) $ & $[4,5; 5) $ \\
			\hline Tần số &  4 & 9 & 14 & 11 & 7 & 5 \\
			\hline 
		\end{tabular}
	\end{center}
	Tính giá trị  đại diện của nhóm $[2{,}5;3)$.
	\loigiai{ 
	Giá trị đại  diện của nhóm $[2{,}5;3)$ là: $\dfrac{2{,}5+3}{2}=2{,}75$.}
\end{ex}

% nguồn: Ok_TuanLuan_Toan11_Lan2
\begin{ex}%[0D6H3-2]%[ma:e522291d]
	Một lần kiểm tra toán của một  lớp gồm $55$ học sinh, thống kê điểm số như sau:
	\begin{center}
		\begin{tabular}{|l|c|c|c|c|c|c|c|c|c|c|c|}
			\hline Điểm &0&1&2&3&4&5&6&7&8&9&10\\ 
			\hline  Số học sinh &0&3&3&5&4&12&10&8&7&1&2\\
			\hline 
		\end{tabular}
	\end{center}
	\begin{enumerate}
		\item Tính số điểm toán trung bình của  $55$ học sinh trên.
		\item Tổng hợp lại số liệu trên vào bảng tần số ghép nhóm theo mẫu  sau:
		\begin{center}
			\begin{tabular}{|l|c|c|c|c|c|}
				\hline Điểm & $[1; 2]  $ & $[3; 4] $ & $[5; 6] $ & $[7; 8] $ & $[9; 10] $ \\
				\hline Số học sinh & $ ? $ & $ ? $ & $ ? $ & $ ? $ & $  ? $ \\
				 \hline
			\end{tabular}
		\end{center}
		\item Hãy  ước lượng trung bình số điểm toán của $55$ học sinh từ bảng tần số ghép nhóm trên.
	\end{enumerate}
	 \loigiai{
		\begin{enumerate}
			 \item Số điểm toán trung bình của $55$ học sinh trên là:\\
			$\overline{x}=(  1\cdot 3+2\cdot 3+3\cdot 5+4\cdot 4+5\cdot 12+6\cdot 10+7\cdot 8+8\cdot 7+9\cdot 1+10\cdot 2): 55\approx 5{,}47$.
			 \item Bảng tần số ghép nhóm:
			\begin{center}
				\begin{tabular}{|l|c|c|c|c|c|}
					\hline Điểm & $[1; 2] $ & $[3; 4] $ & $[5; 6] $ & $[7;  8] $ & $[9; 10] $ \\
					\hline Số học sinh  & 6 & 9 & 22 & 15 & 3 \\
					 \hline
				\end{tabular}
			\end{center}
			\item Hiệu chỉnh  lại bảng số liệu ta có:
			\begin{center}
				\begin{tabular}{|l|c|c|c|c|c|}
					\hline Điểm & $[0,5; 2,5) $ & $[2,5;  4,5) $ & $[4,5; 6,5) $ & $[6,5; 8,5) $ & $[8,5; 10,5) $ \\
					\hline Giá trị đại  diện & 1,5 & 3,5 & 5,5 & 7,5 & 9,5 \\
					\hline Số học sinh & 6 & 9 & 22 & 15  & 3 \\
					\hline 
				\end{tabular}
			\end{center}
			 Trung bình số điểm toán của $55$ học sinh là
			\begin{center}
				$\overline{x}=\dfrac{1{,}5\cdot  6+3{,}5\cdot 9+5{,}5\cdot 22+7{,}5\cdot 15+9{,}5\cdot 3}{55}=5{,}5$.
			\end{center}
		\end{enumerate}
	} 
\end{ex}

% nguồn: Ok_TuanLuan_Toan11_Lan2
\begin{ex}%[0D6H3-2]%[ma:8a881f40]
	Kết quả đo chiều cao của $81$ cây  tre $2$ năm tuổi ở một ngọn đồi được biểu diễn ở biểu đồ dưới đây. Chiều cao $81$ cây tre $2$ năm tuổi.
	\begin{center}
		\begin{tikzpicture}[scale=1,>=stealth,font=\footnotesize,line join=round,line cap=round]
				\draw[->]  (0,0)--(13,0);
				\draw[->] (0,0)--(0,6) node[midway,sloped,above=12pt]{\small{\bf Số  cây}};
				 \foreach \y in {1,2,3,4,5}
				 \pgfmathsetmacro\giatri{int (5*\y)}
				\draw[gray]  (0,\y) node[left]{\color{black}{\giatri}}--(12.5,\y);
				 \foreach \x/\y in {1/1,3/2.2,5/4.6,7/4.8,9/1.8,11/0.8}
				\draw[fill=blue!50] (\x,0)  rectangle +(1,\y);
				\foreach \x/\y/\z/\nhan in {1.5/{$[1; 5]$}/1/{$5$},3.5/{$(5; 9]$}/2.2/{$11$},5.5/{$(9; 13]$}/4.6/{$23$},7.5/{$(13; 17]$}/4.8/{$24$},9.5/{$(17; 21]$}/1.8/{$14$},11.5/{$(21; 25]$}/0.8/{$4$}} 
					 {
				\path (\x,0)  node[below]{\y};
				\draw (\x,\z) node[above]{\nhan}; 
				} 
			\end{tikzpicture}
	\end{center}
	Hãy ước lượng  số trung bình và mốt của mẫu số liệu ghép nhóm trên.
	\loigiai{ 
		\begin{center}
			\begin{tabular}{|l|c|c|c|c|c|c|}
				\hline Chiều cao & $[1; 5] $ & $[5; 9] $ & $[9; 13] $ & $[13;  17] $ & $[17; 21] $ & $[21; 25] $ \\
				\hline Giá trị  đại diện & $3$ & $7$ & $11$ & $15$ & $19$ & $23$ \\
				\hline Số cây  & $5$ & $11$ & $23$ & $24$ & $14$ & $4$ \\
				 \hline
			\end{tabular}
		\end{center}
		Chiều cao trung bình của các cây tre 2 năm tuổi trong ngọn đồi  là:
		\begin{center}
			$\overline{x}=\dfrac{5\cdot 3+11\cdot 7+23\cdot  11+24\cdot 15+14\cdot 19+4\cdot 23}{81}\approx 13{,}123$.
		\end{center}
		Ta có:  ${u_m}=13,{u_{m+1}}=17,{n_m}=24,{n_{m+1}}=14;{n_{m-1}}=23,{u_{m+1}}-{u_m}=4$.\\
		Do  đó mốt của mẫu số liệu ghép nhóm trên là: ${M_0}=13+\dfrac{24-23}{24-23+24-14}\cdot 4\approx 13{,}363$.
	 }
\end{ex}

% nguồn: Ok_TuanLuan_Toan11_Lan2
\begin{ex}%[0D6H3-2]%[ma:d49bab56]
	Người ta đếm số xe ô tô đi qua một trạm thu phí mỗi phút trong  khoảng thời gian từ $11$ giờ đến $11$ giờ $30$ phút buổi trưa. Kết quả được ghi lại ở bảng sau:
	\begin{center}
		\begin{tabular}{|l|l|l|l|l|l|l|l|l|l|l|l|l|l|l|}
			\hline 25 & 21 & 23 & 29 & 37 &  23 & 25 & 31 & 26 & 16 & 32 & 33 & 29 & 16 & 31 \\
			 \hline 25 & 17 & 29 & 21 & 28 & 36 & 24 & 26 & 21 & 18 & 21 & 21 & 19 & 27 & 25 \\
			 \hline
		\end{tabular}
	\end{center}
	\begin{enumerate}
		\item Tính số xe trung bình đi qua  trạm thu phí trong mỗi phút.
		\item Tổng hợp lại số liệu trên vào bảng  tần số ghép nhóm theo mẫu sau:
		\begin{center}
			\begin{tabular}{|l|c|c|c|c|c|}
				\hline Số xe & $[16;  20] $ & $[21; 25] $ & $[26; 30] $ & $[31; 35] $ & $[36; 40] $ \\
				\hline Số lần & $ ? $ & $ ? $ & $ ? $ & $ ? $ & $ ? $ \\ 
				\hline 
			\end{tabular}
		\end{center}
		\item Hãy ước lượng trung bình số xe đi  qua trạm thu phí trong mỗi phút từ bảng tần số ghép nhóm trên.
	\end{enumerate}
	\loigiai{ 
		\begin{enumerate}
			\item Số xe trung bình đi qua trạm  thu phí trong mỗi phút:
			\begin{center}
				$\bar{x}=\dfrac{16\cdot 2+17+18+19+\cdots+32+33+36+37}{30}=25{,}16$. 
			\end{center}
			\item Tổng hợp  lại số liệu trên vào bảng tần số ghép nhóm:
			\begin{center}
				\begin{tabular}{|l|c|c|c|c|c|}
					\hline Số xe & $[16; 20] $ & $[21; 25] $ & $[26; 30] $ & $[31; 35] $ & $[36; 40]  $ \\
					\hline Số lần & 5 & 12  & 7 & 4 & 2 \\
					\hline 
				\end{tabular}
			\end{center}
			\item Ước lượng số xe  trung bình đi qua trạm thu phí trong mỗi phút:
			\begin{center}
				$\bar{x}=\dfrac{18\cdot 5+23\cdot 12+28\cdot 7+33\cdot  4+38\cdot 2}{30}=25{,}66$.
			\end{center}
		\end{enumerate}
	 }
\end{ex}

% nguồn: Ok_TuanLuan_Toan11_Lan2
\begin{ex}%[0D6H3-2]%[ma:4ca61ad4]
	 Một nhà thực vật học đo chiều dài của $74$ chiếc lá và thu được kết quả như sau (đơn vị mm).
	\begin{center}
		\begin{tabular}{|l|c|c|c|c|c|c|}
			\hline Chiều dài & $[5{,}45;  5{,}85) $ & $[5{,}85; 6{,}25) $ & $[6{,}25; 6{,}65) $ & $[6{,}65; 7{,}05) $ & $[7{,}05; 7{,}45) $ & $[7{,}45; 7{,}85) $ \\
			\hline Tần số &  5 & 9 & 15 & 19 & 16 & 8 \\
			 \hline
		\end{tabular}
	\end{center}
	Tính  giá trị đại diện và độ dài của nhóm $[6{,}25;6{,}65)$.
	\loigiai{ 
		Giá trị đại diện của nhóm  là $\dfrac{6{,}25+6{,}65}{2}=6{,}45$.\\
		Độ dài của nhóm là  $6{,}65-6{,}25=0{,}4$.
	 }
\end{ex}

% nguồn: Ok_TuanLuan_Toan11_Lan2
\begin{ex}%[0D6H3-2]%[ma:468490b4]
	Điểm thi  môn Toán (thang điểm 100, điểm được làm tròn đến 1) của $60$ thí sinh được cho trong bảng sau:\\
	\begin{tabular}{|>{\raggedright\arraybackslash}m{1.3cm}|>{\centering\arraybackslash}m{.9cm}|>{\centering\arraybackslash}m{1cm}|>{\centering\arraybackslash}m{1cm}|>{\centering\arraybackslash}m{1cm}|>{\centering\arraybackslash}m{1cm}|>{\centering\arraybackslash}m{1cm}|>{\centering\arraybackslash}m{1cm}|>{\centering\arraybackslash}m{1cm}|>{\centering\arraybackslash}m{1cm}|>{\centering\arraybackslash}m{1.2cm}|}
		\hline Điểm & $[0; 10) $ & $[10; 20) $ & $[20; 30) $ & $[30; 40) $ & $[40; 50) $ & $[50; 60) $ & $[60; 70) $  & $[70; 80) $ & $[80; 90) $ & $[90; 100) $ \\
		\hline 
		 Số thí sinh & 1 & 2 & 4 & 6 & 15 & 12 & 10 & 6 & 3 & 1 \\
		 \hline
	\end{tabular}\\
	Mốt của mẫu số liệu  trên bằng bao nhiêu?
	\loigiai{ 
		Nhóm chứa mốt  của mẫu số liệu là $[40;50)$.\\
		Do đó ${u_m}=40;{n_{m-1}}=6,{n_{m+1}}=12,{u_{m+1}}-{u_m}=50-40=10$.\\ 
		Mốt của  mẫu số liệu ghép nhóm là:
		${M_0}=40+\dfrac{15-6}{( 15-6)+( 15-12)}\cdot 10=47{,}5$. 
	} 
\end{ex}

% nguồn: Ok_TuanLuan_Toan11_Lan2
\begin{ex}%[0D6H3-2]%[ma:d8a5af1f]
	Một Thư viện trường THPT trong tỉnh Quảng Ninh thống kê số lượng sách học sinh được mượn  mỗi ngày trong $60$ ngày ở bảng sau:
	\begin{center}
		\begin{tabular}{|c|c|c|c|c|c|c|c|}
			\hline Số sách & $[16; 20] $ & $[21; 25] $ & $[26; 30] $ & $[31; 35] $ & $[36; 40] $ & $[41; 45] $ & $[46; 50] $  \\
			\hline Số ngày  & 10 & 11 & 14 & 9 & 8 & 5 & 3 \\
			\hline 
		\end{tabular}
	\end{center}
	Hãy  ước lượng số trung bình và mốt của mẫu số liệu ghép nhóm trên.
	 \loigiai{
		Tổng hợp lại  số liệu trên vào bảng tần số ghép nhóm:
		\begin{center}
			\begin{tabular}{|l|c|c|c|c|c|c|c|}
				\hline Số sách & $[16; 20] $ & $[21; 25] $ & $[26; 30] $ & $[31;  35] $ & $[36; 40] $ & $[41; 45] $ & $[46; 50] $ \\
				\hline Giá trị đại diện  & 18 & 23 & 28 & 33 & 38 & 43 & 48 \\
				\hline  Số ngày & 10 & 11 & 14 & 9 & 8 & 5 & 3 \\
				 \hline
			\end{tabular}
		\end{center}
		Số sách trung bình học sinh  mượn trong hai tháng là:
		\begin{center}
			$\bar{x}=\dfrac{18\cdot 10+23\cdot 11+28\cdot 14+33\cdot 9+38\cdot 8+43\cdot  5+48\cdot 3}{60}=29{,}75$.
		\end{center}
		Ta có: ${u_m}=26,{u_{m+1}}=30,{n_m}=14,{n_{m+1}}=9;{n_{m-1}}=11,{u_{m+1}}-{u_m}=4$.\\ 
		Do đó mốt  của mẫu số liệu ghép nhóm trên là: ${M_0}=26+\dfrac{14-11}{14-11+14-9}\cdot 4=27{,}5$.
	 }
\end{ex}

% nguồn: Ok_TuanLuan_Toan11_Lan2
\begin{ex}%[0D6H3-2]%[ma:f70b1899]
	Cân nặng (đơn  vị: kg) của $28$ học sinh lớp 11 được thống kê trong bảng tần số ghép nhóm sau:
	\begin{center}
		\begin{tabular}{|c|c|c|c|c|c|}
			\hline Cân nặng & $[45; 49) $ & $[49; 53) $ & $[53;  57) $ & $[57; 61) $ & $[61; 65) $ \\
			 \hline Số học sinh & 4 & 5 & 7 & 7 & 5 \\
			\hline 
		\end{tabular}
	\end{center}
	 Tính cân nặng trung bình của $28$ học sinh trên.
	 \loigiai{
		Cân nặng trung bình của 28 học  sinh trên xấp xỉ bằng:
		\begin{center}
			$\overline{x}=\dfrac{47\cdot 4+51\cdot 5+55\cdot 7+59\cdot  7+63\cdot 5}{28}\approx 55{,}57$ (kg).
		\end{center}
	 }
\end{ex}

% nguồn: Ok_TuanLuan_Toan11_Lan2
\begin{ex}%[0D6H3-2]%[ma:1eec983d]
	 Quãng đường (km) từ nhà đến nơi làm việc của $40$ công nhân một nhà máy được ghi lại như sau:
	\begin{center}
		\begin{tabular}{cccccccccc}
			5 & 3 & 10 & 20  & 25 & 11 & 13 & 7 & 12 & 31 \\
			19 & 10 & 12 & 17 & 18 & 11 & 32 & 17  & 16 & 2\\
			7 & 9 &  7 & 8 & 3 & 5 & 12 & 15 & 18 & 3 \\
			12 & 14 & 2 &  9 & 6 & 15 & 15 & 7 & 6 & 12
		\end{tabular}
	\end{center}
	Ghép nhóm  số liệu trên thành các nhóm
	$[0;5)$; $[5;10)$; $[10;15)$; $[15;20)$; $[20;25)$; $[25;30)$; $[30;35)$. Tính tần số của nhóm  $[15;20)$.
	\loigiai{ 
	Tần  số của nhóm $[15;20)$ là: $ n=9$.}
\end{ex}

% nguồn: Ok_TuanLuan_Toan11_Lan2
\begin{ex}%[0D6H3-2]%[ma:8d92c473]
	Tổng số lượng mưa trong tháng $8$ đo được tại một trạm quan trắc đặt tại Vũng Tàu từ năm 2002 đến năm  2020 được ghi lại như dưới đây (đơn vị: mm).
	\begin{center}
		\begin{tabular}{ccccccc}
			$121{,}8$&$158{,}3$&$334{,}9$&$200{,}9$&$165{,}6$&$161{,}5$&$194{,}3$\\ 
			$220{,}7$&$189{,}8$&$243{,}2$&$165{,}9$&$165{,}9$&$134$&$173$\\ 
			 $169$&$189$&$254$&$168$&$255$
		\end{tabular}
	\end{center}
	\begin{enumerate}
		\item Hoàn thành bảng tần số  ghép nhóm theo mẫu sau:
		\begin{center}
			\begin{tabular}{|l|c|c|c|c|}
				 \hline
				Tổng lượng mưa trong tháng 8  (mm)&$[120;175)$&$[175;230)$&$[230;285)$&$[285;340)$\\ \hline
				Số  năm&&&&\\ \hline
			\end{tabular}
		\end{center}
		\item Hãy ước lượng lượng mưa trung bình từ mẫu số liệu  ở bảng tần số ghép nhóm trên.
	\end{enumerate}
	 \loigiai{
		\begin{enumerate}
			\item Bảng  tần số ghép nhóm.
			\begin{center}
				\begin{tabular}{|l|c|c|c|c|}
					 \hline
					Tổng lượng mưa  trong tháng 8 (mm)&$[120;175)$&$[175;230)$&$[230;285)$&$[285;340)$\\ \hline
					Số năm&10&5&3&1\\  \hline
				\end{tabular}
			\end{center}
			\item Thống kê lượng  mưa trung bình theo giá trị đại diện.
			\begin{center}
				\begin{tabular}{|l|c|c|c|c|}
					\hline 
					 Tổng lượng mưa trong tháng 8 (mm)&$[120;175)$&$[175;230)$&$[230;285)$&$[285;340)$\\ \hline
					Giá trị đại  diện&$147{,}5$&$202{,}5$&$275{,}5$&$312{,}5$\\ \hline
					Số năm&10&5&3&1\\  \hline
				\end{tabular}
			\end{center}
			Lượng mưa trung bình là $( 147{,}5\cdot 10+202{,}5\cdot 5+257{,}5\cdot 3+312{,}5\cdot 1): 19\approx  188{,}03$.
		\end{enumerate}
	 }
\end{ex}

% nguồn: Ok_TuanLuan_Toan11_Lan2
\begin{ex}%[0D6H3-2]%[ma:cb9de397]
	Thống  kê chiều cao của $60$ học sinh lớp 11 ở một trường trung học phổ thông (đơn vị cm) thu được kết quả như sau:
	\begin{center}
		\begin{tabular}{ccccccccccccc}
			 $175$&$163$&$146$&$150$&$176$&$162$&$147$&$151$&$176$&$161$&$149$&$152$&$177$\\
			$165$&$148$&$153$&$176$&$169$&$152$&$155$&$170$&$144$&$168$&$160$&$170$&$143$\\ 
			$167$&$160$&$170$&$142$&$166$&$160$&$165$&$141$&$174$&$161$&$166$&$144$&$173$\\ 
			$162$&$175$&$156$&$161$&$172$&$175$&$157$&$162$&$171$&$176$&$160$&$158$&$170$\\ 
			 $176$&$164$&$159$&$170$&$175$&$163$&$160$&$170$
		\end{tabular}
	\end{center}
	\begin{enumerate}
		\item Hãy thu gọn dữ liệu bằng cách lập bảng tần số ghép nhóm với các  nhóm ghép là $[135;145)$; $[145;155)$; $[155;165)$; $[165;175)$; $[175;185)$.
		\item Hãy ước lượng chiều cao trung bình của học sinh trong lớp học đó từ bảng tần  số ghép nhóm đã lập ở câu a.
	\end{enumerate}
	 \loigiai{
		\begin{enumerate}
			\item Bảng tần số ghép nhóm chiều cao của $60$  học sinh lớp 11.
			\begin{center}
				\begin{tabular}{|l|c|c|c|c|c|}
					\hline 
					Chiều  cao (cm)&$[135;145)$&$[145;155)$&$[155;165)$&$[165;175)$&$[175;185)$\\
					 \hline
					Tần  số&5&9&19&17&10 \\
					 \hline
				\end{tabular}
			\end{center}
			 \item Ta có bảng thống kê chiều cao theo giá trị đại diện.
			\begin{center}
				\begin{tabular}{|l|c|c|c|c|c|}
					 \hline
					Chiều cao (cm)&$[135;145)$&$[145;155)$&$[155;165)$&$[165;175)$&$[175;185)$\\ 
					\hline 
					Chiều cao đại diện&140&150&160&170&180\\ 
					 \hline
					Tần  số&5&9&19&17&10\\
					\hline 
				\end{tabular}
			\end{center}
			Chiều cao trung bình là  $( 140\cdot 5+150\cdot 9+160\cdot 19+170\cdot 17+180\cdot 10): 60=163$ cm.
	\end{enumerate} }
\end{ex}

% nguồn: Ok_TuanLuan_Toan11_Lan2
\begin{ex}%[0D6H3-2]%[ma:38c95f78]
	Kết quả  kì thi trắc nghiệm môn Toán với thang điểm $100$ của $32$ học sinh được cho trong bảng sau:
	\begin{center}
		\begin{tabular}{|l|c|c|c|c|c|c|}
			\hline Điểm & $[40; 50) $ & $[50; 60) $ & $[60; 70) $ & $[70; 80) $ & $[80;  90) $ & $[90; 100) $ \\
			\hline Tần số & 4 & 6 & 11 & 6 & 3 & 2  \\
			 \hline
		\end{tabular}
	\end{center}
	 Tính độ dài của nhóm $[90;100)$.
	 \loigiai{
		Độ dài của nhóm  $[90;100)$ là $100-90=10$.
	 }
\end{ex}
